\documentclass[11pt]{article}

    \usepackage[breakable]{tcolorbox}
    \usepackage{parskip} % Stop auto-indenting (to mimic markdown behaviour)
    

    % Basic figure setup, for now with no caption control since it's done
    % automatically by Pandoc (which extracts ![](path) syntax from Markdown).
    \usepackage{graphicx}
    % Maintain compatibility with old templates. Remove in nbconvert 6.0
    \let\Oldincludegraphics\includegraphics
    % Ensure that by default, figures have no caption (until we provide a
    % proper Figure object with a Caption API and a way to capture that
    % in the conversion process - todo).
    \usepackage{caption}
    \DeclareCaptionFormat{nocaption}{}
    \captionsetup{format=nocaption,aboveskip=0pt,belowskip=0pt}

    \usepackage{float}
    \floatplacement{figure}{H} % forces figures to be placed at the correct location
    \usepackage{xcolor} % Allow colors to be defined
    \usepackage{enumerate} % Needed for markdown enumerations to work
    \usepackage{geometry} % Used to adjust the document margins
    \usepackage{amsmath} % Equations
    \usepackage{amssymb} % Equations
    \usepackage{textcomp} % defines textquotesingle
    % Hack from http://tex.stackexchange.com/a/47451/13684:
    \AtBeginDocument{%
        \def\PYZsq{\textquotesingle}% Upright quotes in Pygmentized code
    }
    \usepackage{upquote} % Upright quotes for verbatim code
    \usepackage{eurosym} % defines \euro

    \usepackage{iftex}
    \ifPDFTeX
        \usepackage[T1]{fontenc}
        \IfFileExists{alphabeta.sty}{
              \usepackage{alphabeta}
          }{
              \usepackage[mathletters]{ucs}
              \usepackage[utf8x]{inputenc}
          }
    \else
        \usepackage{fontspec}
        \usepackage{unicode-math}
    \fi

    \usepackage{fancyvrb} % verbatim replacement that allows latex
    \usepackage{grffile} % extends the file name processing of package graphics
                         % to support a larger range
    \makeatletter % fix for old versions of grffile with XeLaTeX
    \@ifpackagelater{grffile}{2019/11/01}
    {
      % Do nothing on new versions
    }
    {
      \def\Gread@@xetex#1{%
        \IfFileExists{"\Gin@base".bb}%
        {\Gread@eps{\Gin@base.bb}}%
        {\Gread@@xetex@aux#1}%
      }
    }
    \makeatother
    \usepackage[Export]{adjustbox} % Used to constrain images to a maximum size
    \adjustboxset{max size={0.9\linewidth}{0.9\paperheight}}

    % The hyperref package gives us a pdf with properly built
    % internal navigation ('pdf bookmarks' for the table of contents,
    % internal cross-reference links, web links for URLs, etc.)
    \usepackage{hyperref}
    % The default LaTeX title has an obnoxious amount of whitespace. By default,
    % titling removes some of it. It also provides customization options.
    \usepackage{titling}
    \usepackage{longtable} % longtable support required by pandoc >1.10
    \usepackage{booktabs}  % table support for pandoc > 1.12.2
    \usepackage{array}     % table support for pandoc >= 2.11.3
    \usepackage{calc}      % table minipage width calculation for pandoc >= 2.11.1
    \usepackage[inline]{enumitem} % IRkernel/repr support (it uses the enumerate* environment)
    \usepackage[normalem]{ulem} % ulem is needed to support strikethroughs (\sout)
                                % normalem makes italics be italics, not underlines
    \usepackage{soul}      % strikethrough (\st) support for pandoc >= 3.0.0
    \usepackage{mathrsfs}
    

    
    % Colors for the hyperref package
    \definecolor{urlcolor}{rgb}{0,.145,.698}
    \definecolor{linkcolor}{rgb}{.71,0.21,0.01}
    \definecolor{citecolor}{rgb}{.12,.54,.11}

    % ANSI colors
    \definecolor{ansi-black}{HTML}{3E424D}
    \definecolor{ansi-black-intense}{HTML}{282C36}
    \definecolor{ansi-red}{HTML}{E75C58}
    \definecolor{ansi-red-intense}{HTML}{B22B31}
    \definecolor{ansi-green}{HTML}{00A250}
    \definecolor{ansi-green-intense}{HTML}{007427}
    \definecolor{ansi-yellow}{HTML}{DDB62B}
    \definecolor{ansi-yellow-intense}{HTML}{B27D12}
    \definecolor{ansi-blue}{HTML}{208FFB}
    \definecolor{ansi-blue-intense}{HTML}{0065CA}
    \definecolor{ansi-magenta}{HTML}{D160C4}
    \definecolor{ansi-magenta-intense}{HTML}{A03196}
    \definecolor{ansi-cyan}{HTML}{60C6C8}
    \definecolor{ansi-cyan-intense}{HTML}{258F8F}
    \definecolor{ansi-white}{HTML}{C5C1B4}
    \definecolor{ansi-white-intense}{HTML}{A1A6B2}
    \definecolor{ansi-default-inverse-fg}{HTML}{FFFFFF}
    \definecolor{ansi-default-inverse-bg}{HTML}{000000}

    % common color for the border for error outputs.
    \definecolor{outerrorbackground}{HTML}{FFDFDF}

    % commands and environments needed by pandoc snippets
    % extracted from the output of `pandoc -s`
    \providecommand{\tightlist}{%
      \setlength{\itemsep}{0pt}\setlength{\parskip}{0pt}}
    \DefineVerbatimEnvironment{Highlighting}{Verbatim}{commandchars=\\\{\}}
    % Add ',fontsize=\small' for more characters per line
    \newenvironment{Shaded}{}{}
    \newcommand{\KeywordTok}[1]{\textcolor[rgb]{0.00,0.44,0.13}{\textbf{{#1}}}}
    \newcommand{\DataTypeTok}[1]{\textcolor[rgb]{0.56,0.13,0.00}{{#1}}}
    \newcommand{\DecValTok}[1]{\textcolor[rgb]{0.25,0.63,0.44}{{#1}}}
    \newcommand{\BaseNTok}[1]{\textcolor[rgb]{0.25,0.63,0.44}{{#1}}}
    \newcommand{\FloatTok}[1]{\textcolor[rgb]{0.25,0.63,0.44}{{#1}}}
    \newcommand{\CharTok}[1]{\textcolor[rgb]{0.25,0.44,0.63}{{#1}}}
    \newcommand{\StringTok}[1]{\textcolor[rgb]{0.25,0.44,0.63}{{#1}}}
    \newcommand{\CommentTok}[1]{\textcolor[rgb]{0.38,0.63,0.69}{\textit{{#1}}}}
    \newcommand{\OtherTok}[1]{\textcolor[rgb]{0.00,0.44,0.13}{{#1}}}
    \newcommand{\AlertTok}[1]{\textcolor[rgb]{1.00,0.00,0.00}{\textbf{{#1}}}}
    \newcommand{\FunctionTok}[1]{\textcolor[rgb]{0.02,0.16,0.49}{{#1}}}
    \newcommand{\RegionMarkerTok}[1]{{#1}}
    \newcommand{\ErrorTok}[1]{\textcolor[rgb]{1.00,0.00,0.00}{\textbf{{#1}}}}
    \newcommand{\NormalTok}[1]{{#1}}

    % Additional commands for more recent versions of Pandoc
    \newcommand{\ConstantTok}[1]{\textcolor[rgb]{0.53,0.00,0.00}{{#1}}}
    \newcommand{\SpecialCharTok}[1]{\textcolor[rgb]{0.25,0.44,0.63}{{#1}}}
    \newcommand{\VerbatimStringTok}[1]{\textcolor[rgb]{0.25,0.44,0.63}{{#1}}}
    \newcommand{\SpecialStringTok}[1]{\textcolor[rgb]{0.73,0.40,0.53}{{#1}}}
    \newcommand{\ImportTok}[1]{{#1}}
    \newcommand{\DocumentationTok}[1]{\textcolor[rgb]{0.73,0.13,0.13}{\textit{{#1}}}}
    \newcommand{\AnnotationTok}[1]{\textcolor[rgb]{0.38,0.63,0.69}{\textbf{\textit{{#1}}}}}
    \newcommand{\CommentVarTok}[1]{\textcolor[rgb]{0.38,0.63,0.69}{\textbf{\textit{{#1}}}}}
    \newcommand{\VariableTok}[1]{\textcolor[rgb]{0.10,0.09,0.49}{{#1}}}
    \newcommand{\ControlFlowTok}[1]{\textcolor[rgb]{0.00,0.44,0.13}{\textbf{{#1}}}}
    \newcommand{\OperatorTok}[1]{\textcolor[rgb]{0.40,0.40,0.40}{{#1}}}
    \newcommand{\BuiltInTok}[1]{{#1}}
    \newcommand{\ExtensionTok}[1]{{#1}}
    \newcommand{\PreprocessorTok}[1]{\textcolor[rgb]{0.74,0.48,0.00}{{#1}}}
    \newcommand{\AttributeTok}[1]{\textcolor[rgb]{0.49,0.56,0.16}{{#1}}}
    \newcommand{\InformationTok}[1]{\textcolor[rgb]{0.38,0.63,0.69}{\textbf{\textit{{#1}}}}}
    \newcommand{\WarningTok}[1]{\textcolor[rgb]{0.38,0.63,0.69}{\textbf{\textit{{#1}}}}}


    % Define a nice break command that doesn't care if a line doesn't already
    % exist.
    \def\br{\hspace*{\fill} \\* }
    % Math Jax compatibility definitions
    \def\gt{>}
    \def\lt{<}
    \let\Oldtex\TeX
    \let\Oldlatex\LaTeX
    \renewcommand{\TeX}{\textrm{\Oldtex}}
    \renewcommand{\LaTeX}{\textrm{\Oldlatex}}
    % Document parameters
    % Document title
    \title{tutorial\_notebook}
    
    
    
    
    
    
    
% Pygments definitions
\makeatletter
\def\PY@reset{\let\PY@it=\relax \let\PY@bf=\relax%
    \let\PY@ul=\relax \let\PY@tc=\relax%
    \let\PY@bc=\relax \let\PY@ff=\relax}
\def\PY@tok#1{\csname PY@tok@#1\endcsname}
\def\PY@toks#1+{\ifx\relax#1\empty\else%
    \PY@tok{#1}\expandafter\PY@toks\fi}
\def\PY@do#1{\PY@bc{\PY@tc{\PY@ul{%
    \PY@it{\PY@bf{\PY@ff{#1}}}}}}}
\def\PY#1#2{\PY@reset\PY@toks#1+\relax+\PY@do{#2}}

\@namedef{PY@tok@w}{\def\PY@tc##1{\textcolor[rgb]{0.73,0.73,0.73}{##1}}}
\@namedef{PY@tok@c}{\let\PY@it=\textit\def\PY@tc##1{\textcolor[rgb]{0.24,0.48,0.48}{##1}}}
\@namedef{PY@tok@cp}{\def\PY@tc##1{\textcolor[rgb]{0.61,0.40,0.00}{##1}}}
\@namedef{PY@tok@k}{\let\PY@bf=\textbf\def\PY@tc##1{\textcolor[rgb]{0.00,0.50,0.00}{##1}}}
\@namedef{PY@tok@kp}{\def\PY@tc##1{\textcolor[rgb]{0.00,0.50,0.00}{##1}}}
\@namedef{PY@tok@kt}{\def\PY@tc##1{\textcolor[rgb]{0.69,0.00,0.25}{##1}}}
\@namedef{PY@tok@o}{\def\PY@tc##1{\textcolor[rgb]{0.40,0.40,0.40}{##1}}}
\@namedef{PY@tok@ow}{\let\PY@bf=\textbf\def\PY@tc##1{\textcolor[rgb]{0.67,0.13,1.00}{##1}}}
\@namedef{PY@tok@nb}{\def\PY@tc##1{\textcolor[rgb]{0.00,0.50,0.00}{##1}}}
\@namedef{PY@tok@nf}{\def\PY@tc##1{\textcolor[rgb]{0.00,0.00,1.00}{##1}}}
\@namedef{PY@tok@nc}{\let\PY@bf=\textbf\def\PY@tc##1{\textcolor[rgb]{0.00,0.00,1.00}{##1}}}
\@namedef{PY@tok@nn}{\let\PY@bf=\textbf\def\PY@tc##1{\textcolor[rgb]{0.00,0.00,1.00}{##1}}}
\@namedef{PY@tok@ne}{\let\PY@bf=\textbf\def\PY@tc##1{\textcolor[rgb]{0.80,0.25,0.22}{##1}}}
\@namedef{PY@tok@nv}{\def\PY@tc##1{\textcolor[rgb]{0.10,0.09,0.49}{##1}}}
\@namedef{PY@tok@no}{\def\PY@tc##1{\textcolor[rgb]{0.53,0.00,0.00}{##1}}}
\@namedef{PY@tok@nl}{\def\PY@tc##1{\textcolor[rgb]{0.46,0.46,0.00}{##1}}}
\@namedef{PY@tok@ni}{\let\PY@bf=\textbf\def\PY@tc##1{\textcolor[rgb]{0.44,0.44,0.44}{##1}}}
\@namedef{PY@tok@na}{\def\PY@tc##1{\textcolor[rgb]{0.41,0.47,0.13}{##1}}}
\@namedef{PY@tok@nt}{\let\PY@bf=\textbf\def\PY@tc##1{\textcolor[rgb]{0.00,0.50,0.00}{##1}}}
\@namedef{PY@tok@nd}{\def\PY@tc##1{\textcolor[rgb]{0.67,0.13,1.00}{##1}}}
\@namedef{PY@tok@s}{\def\PY@tc##1{\textcolor[rgb]{0.73,0.13,0.13}{##1}}}
\@namedef{PY@tok@sd}{\let\PY@it=\textit\def\PY@tc##1{\textcolor[rgb]{0.73,0.13,0.13}{##1}}}
\@namedef{PY@tok@si}{\let\PY@bf=\textbf\def\PY@tc##1{\textcolor[rgb]{0.64,0.35,0.47}{##1}}}
\@namedef{PY@tok@se}{\let\PY@bf=\textbf\def\PY@tc##1{\textcolor[rgb]{0.67,0.36,0.12}{##1}}}
\@namedef{PY@tok@sr}{\def\PY@tc##1{\textcolor[rgb]{0.64,0.35,0.47}{##1}}}
\@namedef{PY@tok@ss}{\def\PY@tc##1{\textcolor[rgb]{0.10,0.09,0.49}{##1}}}
\@namedef{PY@tok@sx}{\def\PY@tc##1{\textcolor[rgb]{0.00,0.50,0.00}{##1}}}
\@namedef{PY@tok@m}{\def\PY@tc##1{\textcolor[rgb]{0.40,0.40,0.40}{##1}}}
\@namedef{PY@tok@gh}{\let\PY@bf=\textbf\def\PY@tc##1{\textcolor[rgb]{0.00,0.00,0.50}{##1}}}
\@namedef{PY@tok@gu}{\let\PY@bf=\textbf\def\PY@tc##1{\textcolor[rgb]{0.50,0.00,0.50}{##1}}}
\@namedef{PY@tok@gd}{\def\PY@tc##1{\textcolor[rgb]{0.63,0.00,0.00}{##1}}}
\@namedef{PY@tok@gi}{\def\PY@tc##1{\textcolor[rgb]{0.00,0.52,0.00}{##1}}}
\@namedef{PY@tok@gr}{\def\PY@tc##1{\textcolor[rgb]{0.89,0.00,0.00}{##1}}}
\@namedef{PY@tok@ge}{\let\PY@it=\textit}
\@namedef{PY@tok@gs}{\let\PY@bf=\textbf}
\@namedef{PY@tok@gp}{\let\PY@bf=\textbf\def\PY@tc##1{\textcolor[rgb]{0.00,0.00,0.50}{##1}}}
\@namedef{PY@tok@go}{\def\PY@tc##1{\textcolor[rgb]{0.44,0.44,0.44}{##1}}}
\@namedef{PY@tok@gt}{\def\PY@tc##1{\textcolor[rgb]{0.00,0.27,0.87}{##1}}}
\@namedef{PY@tok@err}{\def\PY@bc##1{{\setlength{\fboxsep}{\string -\fboxrule}\fcolorbox[rgb]{1.00,0.00,0.00}{1,1,1}{\strut ##1}}}}
\@namedef{PY@tok@kc}{\let\PY@bf=\textbf\def\PY@tc##1{\textcolor[rgb]{0.00,0.50,0.00}{##1}}}
\@namedef{PY@tok@kd}{\let\PY@bf=\textbf\def\PY@tc##1{\textcolor[rgb]{0.00,0.50,0.00}{##1}}}
\@namedef{PY@tok@kn}{\let\PY@bf=\textbf\def\PY@tc##1{\textcolor[rgb]{0.00,0.50,0.00}{##1}}}
\@namedef{PY@tok@kr}{\let\PY@bf=\textbf\def\PY@tc##1{\textcolor[rgb]{0.00,0.50,0.00}{##1}}}
\@namedef{PY@tok@bp}{\def\PY@tc##1{\textcolor[rgb]{0.00,0.50,0.00}{##1}}}
\@namedef{PY@tok@fm}{\def\PY@tc##1{\textcolor[rgb]{0.00,0.00,1.00}{##1}}}
\@namedef{PY@tok@vc}{\def\PY@tc##1{\textcolor[rgb]{0.10,0.09,0.49}{##1}}}
\@namedef{PY@tok@vg}{\def\PY@tc##1{\textcolor[rgb]{0.10,0.09,0.49}{##1}}}
\@namedef{PY@tok@vi}{\def\PY@tc##1{\textcolor[rgb]{0.10,0.09,0.49}{##1}}}
\@namedef{PY@tok@vm}{\def\PY@tc##1{\textcolor[rgb]{0.10,0.09,0.49}{##1}}}
\@namedef{PY@tok@sa}{\def\PY@tc##1{\textcolor[rgb]{0.73,0.13,0.13}{##1}}}
\@namedef{PY@tok@sb}{\def\PY@tc##1{\textcolor[rgb]{0.73,0.13,0.13}{##1}}}
\@namedef{PY@tok@sc}{\def\PY@tc##1{\textcolor[rgb]{0.73,0.13,0.13}{##1}}}
\@namedef{PY@tok@dl}{\def\PY@tc##1{\textcolor[rgb]{0.73,0.13,0.13}{##1}}}
\@namedef{PY@tok@s2}{\def\PY@tc##1{\textcolor[rgb]{0.73,0.13,0.13}{##1}}}
\@namedef{PY@tok@sh}{\def\PY@tc##1{\textcolor[rgb]{0.73,0.13,0.13}{##1}}}
\@namedef{PY@tok@s1}{\def\PY@tc##1{\textcolor[rgb]{0.73,0.13,0.13}{##1}}}
\@namedef{PY@tok@mb}{\def\PY@tc##1{\textcolor[rgb]{0.40,0.40,0.40}{##1}}}
\@namedef{PY@tok@mf}{\def\PY@tc##1{\textcolor[rgb]{0.40,0.40,0.40}{##1}}}
\@namedef{PY@tok@mh}{\def\PY@tc##1{\textcolor[rgb]{0.40,0.40,0.40}{##1}}}
\@namedef{PY@tok@mi}{\def\PY@tc##1{\textcolor[rgb]{0.40,0.40,0.40}{##1}}}
\@namedef{PY@tok@il}{\def\PY@tc##1{\textcolor[rgb]{0.40,0.40,0.40}{##1}}}
\@namedef{PY@tok@mo}{\def\PY@tc##1{\textcolor[rgb]{0.40,0.40,0.40}{##1}}}
\@namedef{PY@tok@ch}{\let\PY@it=\textit\def\PY@tc##1{\textcolor[rgb]{0.24,0.48,0.48}{##1}}}
\@namedef{PY@tok@cm}{\let\PY@it=\textit\def\PY@tc##1{\textcolor[rgb]{0.24,0.48,0.48}{##1}}}
\@namedef{PY@tok@cpf}{\let\PY@it=\textit\def\PY@tc##1{\textcolor[rgb]{0.24,0.48,0.48}{##1}}}
\@namedef{PY@tok@c1}{\let\PY@it=\textit\def\PY@tc##1{\textcolor[rgb]{0.24,0.48,0.48}{##1}}}
\@namedef{PY@tok@cs}{\let\PY@it=\textit\def\PY@tc##1{\textcolor[rgb]{0.24,0.48,0.48}{##1}}}

\def\PYZbs{\char`\\}
\def\PYZus{\char`\_}
\def\PYZob{\char`\{}
\def\PYZcb{\char`\}}
\def\PYZca{\char`\^}
\def\PYZam{\char`\&}
\def\PYZlt{\char`\<}
\def\PYZgt{\char`\>}
\def\PYZsh{\char`\#}
\def\PYZpc{\char`\%}
\def\PYZdl{\char`\$}
\def\PYZhy{\char`\-}
\def\PYZsq{\char`\'}
\def\PYZdq{\char`\"}
\def\PYZti{\char`\~}
% for compatibility with earlier versions
\def\PYZat{@}
\def\PYZlb{[}
\def\PYZrb{]}
\makeatother


    % For linebreaks inside Verbatim environment from package fancyvrb.
    \makeatletter
        \newbox\Wrappedcontinuationbox
        \newbox\Wrappedvisiblespacebox
        \newcommand*\Wrappedvisiblespace {\textcolor{red}{\textvisiblespace}}
        \newcommand*\Wrappedcontinuationsymbol {\textcolor{red}{\llap{\tiny$\m@th\hookrightarrow$}}}
        \newcommand*\Wrappedcontinuationindent {3ex }
        \newcommand*\Wrappedafterbreak {\kern\Wrappedcontinuationindent\copy\Wrappedcontinuationbox}
        % Take advantage of the already applied Pygments mark-up to insert
        % potential linebreaks for TeX processing.
        %        {, <, #, %, $, ' and ": go to next line.
        %        _, }, ^, &, >, - and ~: stay at end of broken line.
        % Use of \textquotesingle for straight quote.
        \newcommand*\Wrappedbreaksatspecials {%
            \def\PYGZus{\discretionary{\char`\_}{\Wrappedafterbreak}{\char`\_}}%
            \def\PYGZob{\discretionary{}{\Wrappedafterbreak\char`\{}{\char`\{}}%
            \def\PYGZcb{\discretionary{\char`\}}{\Wrappedafterbreak}{\char`\}}}%
            \def\PYGZca{\discretionary{\char`\^}{\Wrappedafterbreak}{\char`\^}}%
            \def\PYGZam{\discretionary{\char`\&}{\Wrappedafterbreak}{\char`\&}}%
            \def\PYGZlt{\discretionary{}{\Wrappedafterbreak\char`\<}{\char`\<}}%
            \def\PYGZgt{\discretionary{\char`\>}{\Wrappedafterbreak}{\char`\>}}%
            \def\PYGZsh{\discretionary{}{\Wrappedafterbreak\char`\#}{\char`\#}}%
            \def\PYGZpc{\discretionary{}{\Wrappedafterbreak\char`\%}{\char`\%}}%
            \def\PYGZdl{\discretionary{}{\Wrappedafterbreak\char`\$}{\char`\$}}%
            \def\PYGZhy{\discretionary{\char`\-}{\Wrappedafterbreak}{\char`\-}}%
            \def\PYGZsq{\discretionary{}{\Wrappedafterbreak\textquotesingle}{\textquotesingle}}%
            \def\PYGZdq{\discretionary{}{\Wrappedafterbreak\char`\"}{\char`\"}}%
            \def\PYGZti{\discretionary{\char`\~}{\Wrappedafterbreak}{\char`\~}}%
        }
        % Some characters . , ; ? ! / are not pygmentized.
        % This macro makes them "active" and they will insert potential linebreaks
        \newcommand*\Wrappedbreaksatpunct {%
            \lccode`\~`\.\lowercase{\def~}{\discretionary{\hbox{\char`\.}}{\Wrappedafterbreak}{\hbox{\char`\.}}}%
            \lccode`\~`\,\lowercase{\def~}{\discretionary{\hbox{\char`\,}}{\Wrappedafterbreak}{\hbox{\char`\,}}}%
            \lccode`\~`\;\lowercase{\def~}{\discretionary{\hbox{\char`\;}}{\Wrappedafterbreak}{\hbox{\char`\;}}}%
            \lccode`\~`\:\lowercase{\def~}{\discretionary{\hbox{\char`\:}}{\Wrappedafterbreak}{\hbox{\char`\:}}}%
            \lccode`\~`\?\lowercase{\def~}{\discretionary{\hbox{\char`\?}}{\Wrappedafterbreak}{\hbox{\char`\?}}}%
            \lccode`\~`\!\lowercase{\def~}{\discretionary{\hbox{\char`\!}}{\Wrappedafterbreak}{\hbox{\char`\!}}}%
            \lccode`\~`\/\lowercase{\def~}{\discretionary{\hbox{\char`\/}}{\Wrappedafterbreak}{\hbox{\char`\/}}}%
            \catcode`\.\active
            \catcode`\,\active
            \catcode`\;\active
            \catcode`\:\active
            \catcode`\?\active
            \catcode`\!\active
            \catcode`\/\active
            \lccode`\~`\~
        }
    \makeatother

    \let\OriginalVerbatim=\Verbatim
    \makeatletter
    \renewcommand{\Verbatim}[1][1]{%
        %\parskip\z@skip
        \sbox\Wrappedcontinuationbox {\Wrappedcontinuationsymbol}%
        \sbox\Wrappedvisiblespacebox {\FV@SetupFont\Wrappedvisiblespace}%
        \def\FancyVerbFormatLine ##1{\hsize\linewidth
            \vtop{\raggedright\hyphenpenalty\z@\exhyphenpenalty\z@
                \doublehyphendemerits\z@\finalhyphendemerits\z@
                \strut ##1\strut}%
        }%
        % If the linebreak is at a space, the latter will be displayed as visible
        % space at end of first line, and a continuation symbol starts next line.
        % Stretch/shrink are however usually zero for typewriter font.
        \def\FV@Space {%
            \nobreak\hskip\z@ plus\fontdimen3\font minus\fontdimen4\font
            \discretionary{\copy\Wrappedvisiblespacebox}{\Wrappedafterbreak}
            {\kern\fontdimen2\font}%
        }%

        % Allow breaks at special characters using \PYG... macros.
        \Wrappedbreaksatspecials
        % Breaks at punctuation characters . , ; ? ! and / need catcode=\active
        \OriginalVerbatim[#1,codes*=\Wrappedbreaksatpunct]%
    }
    \makeatother

    % Exact colors from NB
    \definecolor{incolor}{HTML}{303F9F}
    \definecolor{outcolor}{HTML}{D84315}
    \definecolor{cellborder}{HTML}{CFCFCF}
    \definecolor{cellbackground}{HTML}{F7F7F7}

    % prompt
    \makeatletter
    \newcommand{\boxspacing}{\kern\kvtcb@left@rule\kern\kvtcb@boxsep}
    \makeatother
    \newcommand{\prompt}[4]{
        {\ttfamily\llap{{\color{#2}[#3]:\hspace{3pt}#4}}\vspace{-\baselineskip}}
    }
    

    
    % Prevent overflowing lines due to hard-to-break entities
    \sloppy
    % Setup hyperref package
    \hypersetup{
      breaklinks=true,  % so long urls are correctly broken across lines
      colorlinks=true,
      urlcolor=urlcolor,
      linkcolor=linkcolor,
      citecolor=citecolor,
      }
    % Slightly bigger margins than the latex defaults
    
    \geometry{verbose,tmargin=1in,bmargin=1in,lmargin=1in,rmargin=1in}
    
    

\begin{document}
    
    \maketitle
    
    

    
    \hypertarget{clsbiweaklens-tutorial}{%
\section{CLSBIWeakLens Tutorial}\label{clsbiweaklens-tutorial}}

We are going to walk through a simple example of performing inference
with CLSBIWeakLens.

First, let's make sure that we have all of our dependencies ready.

    \begin{tcolorbox}[breakable, size=fbox, boxrule=1pt, pad at break*=1mm,colback=cellbackground, colframe=cellborder]
\prompt{In}{incolor}{110}{\boxspacing}
\begin{Verbatim}[commandchars=\\\{\}]
\PY{c+c1}{\PYZsh{} Update PIP}
\PY{o}{\PYZpc{}}\PY{k}{pip} install \PYZhy{}\PYZhy{}upgrade pip \PYZhy{}q

\PY{c+c1}{\PYZsh{} Cosmology}
\PY{o}{\PYZpc{}}\PY{k}{pip} install colossus \PYZhy{}q

\PY{c+c1}{\PYZsh{} Inference}
\PY{o}{\PYZpc{}}\PY{k}{pip} install sbi \PYZhy{}q
\PY{o}{\PYZpc{}}\PY{k}{pip} install emcee \PYZhy{}q

\PY{c+c1}{\PYZsh{} Plotting}
\PY{o}{\PYZpc{}}\PY{k}{pip} install pygtc \PYZhy{}q
\PY{o}{\PYZpc{}}\PY{k}{pip} install chainconsumer \PYZhy{}q
\PY{o}{\PYZpc{}}\PY{k}{pip} install \PYZhy{}\PYZhy{}upgrade ipython \PYZhy{}q
\end{Verbatim}
\end{tcolorbox}

    \begin{Verbatim}[commandchars=\\\{\}]
Note: you may need to restart the kernel to use updated packages.
Note: you may need to restart the kernel to use updated packages.
Note: you may need to restart the kernel to use updated packages.
Note: you may need to restart the kernel to use updated packages.
Note: you may need to restart the kernel to use updated packages.
Note: you may need to restart the kernel to use updated packages.
Note: you may need to restart the kernel to use updated packages.
    \end{Verbatim}

    \begin{tcolorbox}[breakable, size=fbox, boxrule=1pt, pad at break*=1mm,colback=cellbackground, colframe=cellborder]
\prompt{In}{incolor}{89}{\boxspacing}
\begin{Verbatim}[commandchars=\\\{\}]
\PY{k+kn}{import} \PY{n+nn}{numpy} \PY{k}{as} \PY{n+nn}{np}
\PY{k+kn}{import} \PY{n+nn}{json}
\PY{k+kn}{from} \PY{n+nn}{matplotlib} \PY{k+kn}{import} \PY{n}{pyplot} \PY{k}{as} \PY{n}{plt}
\PY{k+kn}{import} \PY{n+nn}{matplotlib}\PY{n+nn}{.}\PY{n+nn}{image} \PY{k}{as} \PY{n+nn}{mpimg}
\end{Verbatim}
\end{tcolorbox}

    \hypertarget{simulation}{%
\section{Simulation}\label{simulation}}

To train our inferrer for Simulation-Based Inference (SBI), we'll need
to generate some simulations. These are the data points from which we
can learn the relationship between NFW profiles (observables) and
log10mass-concentration pairs (parameters) that we'll use later in this
tutorial for inference.

\hypertarget{configuration}{%
\subsection{Configuration}\label{configuration}}

First, we'll load in an example simulation config and look at the
example config values. For the different numerical experiments, we'll
tweak these and use different combinations of configs, but this is a
good representative example of what a sim config could be.

    \begin{tcolorbox}[breakable, size=fbox, boxrule=1pt, pad at break*=1mm,colback=cellbackground, colframe=cellborder]
\prompt{In}{incolor}{4}{\boxspacing}
\begin{Verbatim}[commandchars=\\\{\}]
\PY{c+c1}{\PYZsh{} Command line arguments for generating simulations}
\PY{n}{SIM\PYZus{}ID} \PY{o}{=} \PY{l+s+s1}{\PYZsq{}}\PY{l+s+s1}{sim\PYZus{}z1}\PY{l+s+s1}{\PYZsq{}}
\PY{n}{NUM\PYZus{}SIMS} \PY{o}{=} \PY{l+m+mi}{10000}

\PY{k}{with} \PY{n+nb}{open}\PY{p}{(}\PY{l+s+sa}{f}\PY{l+s+s1}{\PYZsq{}}\PY{l+s+s1}{configs/simulations/}\PY{l+s+si}{\PYZob{}}\PY{n}{SIM\PYZus{}ID}\PY{l+s+si}{\PYZcb{}}\PY{l+s+s1}{.json}\PY{l+s+s1}{\PYZsq{}}\PY{p}{,} \PY{l+s+s1}{\PYZsq{}}\PY{l+s+s1}{r}\PY{l+s+s1}{\PYZsq{}}\PY{p}{)} \PY{k}{as} \PY{n}{f}\PY{p}{:}
    \PY{n}{sim\PYZus{}config} \PY{o}{=} \PY{n}{json}\PY{o}{.}\PY{n}{load}\PY{p}{(}\PY{n}{f}\PY{p}{)}
    
\PY{n+nb}{print}\PY{p}{(}\PY{n}{json}\PY{o}{.}\PY{n}{dumps}\PY{p}{(}\PY{n}{sim\PYZus{}config}\PY{p}{,} \PY{n}{indent}\PY{o}{=}\PY{l+m+mi}{2}\PY{p}{)}\PY{p}{)}
\end{Verbatim}
\end{tcolorbox}

    \begin{Verbatim}[commandchars=\\\{\}]
\{
  "mc\_relation": "child18",
  "mc\_scatter": 0.1,
  "min\_log10mass": 13,
  "max\_log10mass": 16,
  "num\_radial\_bins": 30,
  "profile\_noise\_dex": 0.3,
  "mc\_pair\_subselect": "all",
  "min\_z": 0.2,
  "max\_z": 0.35
\}
    \end{Verbatim}

    \hypertarget{generating-simulations}{%
\subsection{Generating Simulations}\label{generating-simulations}}

To use this sim\_config to generate our simulations, we do the
following: 1) randomly select \texttt{NUM\_SIMS} log10masses in the
defined range (\texttt{min\_log10mass}, \texttt{max\_log10mass}) 2) use
the \texttt{mc\_relation} (with some \texttt{mc\_scatter} and redshift,
\texttt{z}) to find the corresponding concentrations 3) filter these
mc\_pairs based on the \texttt{mc\_pair\_subselect} criteria 4) generate
NFW profiles with \texttt{num\_radial\_bins} radial bins and some
\texttt{profile\_noise\_dex} applied to them 6) output the
\texttt{sample\_mc\_pairs} and the \texttt{simulated\_nfw\_profiles} as
\texttt{.npy} files

    \begin{tcolorbox}[breakable, size=fbox, boxrule=1pt, pad at break*=1mm,colback=cellbackground, colframe=cellborder]
\prompt{In}{incolor}{5}{\boxspacing}
\begin{Verbatim}[commandchars=\\\{\}]
\PY{o}{\PYZpc{}}\PY{k}{run} scripts/gen\PYZus{}simulations.py \PYZhy{}\PYZhy{}sim\PYZus{}id \PYZob{}SIM\PYZus{}ID\PYZcb{} \PYZhy{}\PYZhy{}num\PYZus{}sims \PYZob{}NUM\PYZus{}SIMS\PYZcb{} \PYZhy{}\PYZhy{}regenerate
\end{Verbatim}
\end{tcolorbox}

    \begin{Verbatim}[commandchars=\\\{\}]
Overwriting existing simulation because of --regenerate flag
    \end{Verbatim}

    \begin{tcolorbox}[breakable, size=fbox, boxrule=1pt, pad at break*=1mm,colback=cellbackground, colframe=cellborder]
\prompt{In}{incolor}{6}{\boxspacing}
\begin{Verbatim}[commandchars=\\\{\}]
\PY{c+c1}{\PYZsh{} Confirming the output files were generated correctly}
\PY{n}{sim\PYZus{}output\PYZus{}dir} \PY{o}{=} \PY{l+s+sa}{f}\PY{l+s+s1}{\PYZsq{}}\PY{l+s+s1}{outputs/simulations/}\PY{l+s+si}{\PYZob{}}\PY{n}{SIM\PYZus{}ID}\PY{l+s+si}{\PYZcb{}}\PY{l+s+s1}{.}\PY{l+s+si}{\PYZob{}}\PY{n}{NUM\PYZus{}SIMS}\PY{l+s+si}{\PYZcb{}}\PY{l+s+s1}{\PYZsq{}}
\PY{o}{\PYZpc{}}\PY{k}{ls} \PYZob{}sim\PYZus{}output\PYZus{}dir\PYZcb{}
\end{Verbatim}
\end{tcolorbox}

    \begin{Verbatim}[commandchars=\\\{\}]
sample\_mc\_pairs.npy         simulated\_nfw\_profiles.npy
    \end{Verbatim}

    Let's look at the \texttt{sample\_mc\_pairs} to make sure they look
correct. We expect to see our predicted log10mass-concentration relation
with some scatter applied to it.

    \begin{tcolorbox}[breakable, size=fbox, boxrule=1pt, pad at break*=1mm,colback=cellbackground, colframe=cellborder]
\prompt{In}{incolor}{7}{\boxspacing}
\begin{Verbatim}[commandchars=\\\{\}]
\PY{n}{sample\PYZus{}mc\PYZus{}pairs} \PY{o}{=} \PY{n}{np}\PY{o}{.}\PY{n}{load}\PY{p}{(}\PY{l+s+sa}{f}\PY{l+s+s1}{\PYZsq{}}\PY{l+s+si}{\PYZob{}}\PY{n}{sim\PYZus{}output\PYZus{}dir}\PY{l+s+si}{\PYZcb{}}\PY{l+s+s1}{/sample\PYZus{}mc\PYZus{}pairs.npy}\PY{l+s+s1}{\PYZsq{}}\PY{p}{)}

\PY{c+c1}{\PYZsh{} Expected shape: (NUM\PYZus{}SIMS, 2)}
\PY{n+nb}{print}\PY{p}{(}\PY{n}{np}\PY{o}{.}\PY{n}{shape}\PY{p}{(}\PY{n}{sample\PYZus{}mc\PYZus{}pairs}\PY{p}{)}\PY{p}{)}

\PY{c+c1}{\PYZsh{} 5 example mc\PYZus{}pairs}
\PY{n+nb}{print}\PY{p}{(}\PY{n}{sample\PYZus{}mc\PYZus{}pairs}\PY{p}{[}\PY{p}{:}\PY{l+m+mi}{5}\PY{p}{]}\PY{p}{)}

\PY{c+c1}{\PYZsh{} mc\PYZus{}pairs scatter plot}
\PY{n}{plt}\PY{o}{.}\PY{n}{scatter}\PY{p}{(}\PY{n}{sample\PYZus{}mc\PYZus{}pairs}\PY{p}{[}\PY{p}{:}\PY{p}{,}\PY{l+m+mi}{0}\PY{p}{]}\PY{p}{,} \PY{n}{sample\PYZus{}mc\PYZus{}pairs}\PY{p}{[}\PY{p}{:}\PY{p}{,}\PY{l+m+mi}{1}\PY{p}{]}\PY{p}{,}\PY{n}{s}\PY{o}{=}\PY{l+m+mi}{5}\PY{p}{)}
\PY{n}{plt}\PY{o}{.}\PY{n}{xlabel}\PY{p}{(}\PY{l+s+s1}{\PYZsq{}}\PY{l+s+s1}{log\PYZdl{}\PYZus{}}\PY{l+s+si}{\PYZob{}10\PYZcb{}}\PY{l+s+s1}{\PYZdl{}M [M\PYZdl{}\PYZus{}}\PY{l+s+s1}{\PYZbs{}}\PY{l+s+s1}{odot\PYZdl{}]}\PY{l+s+s1}{\PYZsq{}}\PY{p}{,} \PY{n}{fontsize}\PY{o}{=}\PY{l+s+s1}{\PYZsq{}}\PY{l+s+s1}{xx\PYZhy{}large}\PY{l+s+s1}{\PYZsq{}}\PY{p}{)}
\PY{n}{plt}\PY{o}{.}\PY{n}{ylabel}\PY{p}{(}\PY{l+s+s1}{\PYZsq{}}\PY{l+s+s1}{Concentration}\PY{l+s+s1}{\PYZsq{}}\PY{p}{,} \PY{n}{fontsize}\PY{o}{=}\PY{l+s+s1}{\PYZsq{}}\PY{l+s+s1}{xx\PYZhy{}large}\PY{l+s+s1}{\PYZsq{}}\PY{p}{)}
\PY{n}{plt}\PY{o}{.}\PY{n}{title}\PY{p}{(}\PY{l+s+s1}{\PYZsq{}}\PY{l+s+s1}{Simulated mc\PYZus{}pairs}\PY{l+s+s1}{\PYZsq{}}\PY{p}{)}
\PY{n}{plt}\PY{o}{.}\PY{n}{show}\PY{p}{(}\PY{p}{)}
\end{Verbatim}
\end{tcolorbox}

    \begin{Verbatim}[commandchars=\\\{\}]
(10000, 2)
[[14.93996807  4.32874451]
 [13.86250645  5.28729269]
 [13.24259073  6.07611466]
 [13.67707303  5.55403392]
 [15.93422363  3.88254974]]
    \end{Verbatim}

    \begin{center}
    \adjustimage{max size={0.9\linewidth}{0.9\paperheight}}{tutorial_notebook_files/tutorial_notebook_9_1.png}
    \end{center}
    { \hspace*{\fill} \\}
    
    \hypertarget{observations}{%
\section{Observations}\label{observations}}

Now, let's generate some ``observations'', which are mock NFW profiles
that we'll treat as empirical data. We'll use these for both MCMC and
SBI later in the notebook.

\hypertarget{configuration}{%
\subsection{Configuration}\label{configuration}}

Similar to our approach above, let's look at an example observation
config so we can inspect what we can tweak in numerical experiments.

    \begin{tcolorbox}[breakable, size=fbox, boxrule=1pt, pad at break*=1mm,colback=cellbackground, colframe=cellborder]
\prompt{In}{incolor}{8}{\boxspacing}
\begin{Verbatim}[commandchars=\\\{\}]
\PY{c+c1}{\PYZsh{} Command line arguments for generating simulations}
\PY{n}{OBS\PYZus{}ID} \PY{o}{=} \PY{l+s+s1}{\PYZsq{}}\PY{l+s+s1}{obs\PYZus{}z1\PYZus{}lambda5}\PY{l+s+s1}{\PYZsq{}}
\PY{n}{NUM\PYZus{}OBS} \PY{o}{=} \PY{l+m+mi}{10}

\PY{k}{with} \PY{n+nb}{open}\PY{p}{(}\PY{l+s+sa}{f}\PY{l+s+s1}{\PYZsq{}}\PY{l+s+s1}{configs/observations/}\PY{l+s+si}{\PYZob{}}\PY{n}{OBS\PYZus{}ID}\PY{l+s+si}{\PYZcb{}}\PY{l+s+s1}{.json}\PY{l+s+s1}{\PYZsq{}}\PY{p}{,} \PY{l+s+s1}{\PYZsq{}}\PY{l+s+s1}{r}\PY{l+s+s1}{\PYZsq{}}\PY{p}{)} \PY{k}{as} \PY{n}{f}\PY{p}{:}
    \PY{n}{obs\PYZus{}config} \PY{o}{=} \PY{n}{json}\PY{o}{.}\PY{n}{load}\PY{p}{(}\PY{n}{f}\PY{p}{)}
    
\PY{n+nb}{print}\PY{p}{(}\PY{n}{json}\PY{o}{.}\PY{n}{dumps}\PY{p}{(}\PY{n}{obs\PYZus{}config}\PY{p}{,} \PY{n}{indent}\PY{o}{=}\PY{l+m+mi}{2}\PY{p}{)}\PY{p}{)}
\end{Verbatim}
\end{tcolorbox}

    \begin{Verbatim}[commandchars=\\\{\}]
\{
  "rm\_relation": "mcclintock18",
  "mc\_relation": "child18",
  "min\_richness": 30,
  "max\_richness": 45,
  "rm\_scatter": 0.1,
  "mc\_scatter": 0.1,
  "profile\_noise\_dex": 0.3,
  "num\_radial\_bins": 30,
  "min\_z": 0.2,
  "max\_z": 0.35
\}
    \end{Verbatim}

    \hypertarget{generating-observations}{%
\subsection{Generating Observations}\label{generating-observations}}

To use this obs\_config to generate our observations, we do the
following: 1) randomly select \texttt{NUM\_OBS} richnesses (\(\lambda\))
in the defined range (\texttt{min\_richness}, \texttt{max\_richness}) 2)
use the \texttt{rm\_relation} (with some \texttt{rm\_scatter}) to find
the corresponding log10masses 3) use the \texttt{mc\_relation} (with
some \texttt{mc\_scatter} and redshift, \texttt{z}) to find the
corresponding concentrations 4) generate NFW profiles with
\texttt{num\_radial\_bins} radial bins and some
\texttt{profile\_noise\_dex} applied to them 5) output the
\texttt{drawn\_mc\_pairs} and the \texttt{drawn\_nfw\_profiles} as
\texttt{.npy} files

    \begin{tcolorbox}[breakable, size=fbox, boxrule=1pt, pad at break*=1mm,colback=cellbackground, colframe=cellborder]
\prompt{In}{incolor}{10}{\boxspacing}
\begin{Verbatim}[commandchars=\\\{\}]
\PY{o}{\PYZpc{}}\PY{k}{run} scripts/gen\PYZus{}observations.py \PYZhy{}\PYZhy{}obs\PYZus{}id \PYZob{}OBS\PYZus{}ID\PYZcb{} \PYZhy{}\PYZhy{}num\PYZus{}obs \PYZob{}NUM\PYZus{}OBS\PYZcb{} \PYZhy{}\PYZhy{}regenerate
\end{Verbatim}
\end{tcolorbox}

    \begin{Verbatim}[commandchars=\\\{\}]
Overwriting existing observations because of --regenerate flag
    \end{Verbatim}

    \begin{tcolorbox}[breakable, size=fbox, boxrule=1pt, pad at break*=1mm,colback=cellbackground, colframe=cellborder]
\prompt{In}{incolor}{11}{\boxspacing}
\begin{Verbatim}[commandchars=\\\{\}]
\PY{c+c1}{\PYZsh{} Confirming the output files were generated correctly}
\PY{n}{obs\PYZus{}output\PYZus{}dir} \PY{o}{=} \PY{l+s+sa}{f}\PY{l+s+s1}{\PYZsq{}}\PY{l+s+s1}{outputs/observations/}\PY{l+s+si}{\PYZob{}}\PY{n}{OBS\PYZus{}ID}\PY{l+s+si}{\PYZcb{}}\PY{l+s+s1}{.}\PY{l+s+si}{\PYZob{}}\PY{n}{NUM\PYZus{}OBS}\PY{l+s+si}{\PYZcb{}}\PY{l+s+s1}{\PYZsq{}}
\PY{o}{\PYZpc{}}\PY{k}{ls} \PYZob{}obs\PYZus{}output\PYZus{}dir\PYZcb{}
\end{Verbatim}
\end{tcolorbox}

    \begin{Verbatim}[commandchars=\\\{\}]
drawn\_mc\_pairs.npy                log\_sigmas.npy
drawn\_mc\_pairs.pdf                noiseless\_drawn\_nfw\_profiles.npy
drawn\_mc\_pairs.png                noiseless\_drawn\_nfw\_profiles.pdf
drawn\_nfw\_profiles.npy            noiseless\_drawn\_nfw\_profiles.png
drawn\_nfw\_profiles.pdf            sigmas.npy
drawn\_nfw\_profiles.png
    \end{Verbatim}

    First, let's look at the drawn\_mc\_pairs from the given richness bin.

Note that we generate these plots (\texttt{drawn\_mc\_pairs.png},
\texttt{drawn\_nfw\_profiles.png}) automatically using
\texttt{plotutils.py}, but it's helpful to replot manually here to build
familiarity with the core data objects.

    \begin{tcolorbox}[breakable, size=fbox, boxrule=1pt, pad at break*=1mm,colback=cellbackground, colframe=cellborder]
\prompt{In}{incolor}{112}{\boxspacing}
\begin{Verbatim}[commandchars=\\\{\}]
\PY{n}{drawn\PYZus{}mc\PYZus{}pairs} \PY{o}{=} \PY{n}{np}\PY{o}{.}\PY{n}{load}\PY{p}{(}\PY{l+s+sa}{f}\PY{l+s+s1}{\PYZsq{}}\PY{l+s+si}{\PYZob{}}\PY{n}{obs\PYZus{}output\PYZus{}dir}\PY{l+s+si}{\PYZcb{}}\PY{l+s+s1}{/drawn\PYZus{}mc\PYZus{}pairs.npy}\PY{l+s+s1}{\PYZsq{}}\PY{p}{)}
\PY{n}{drawn\PYZus{}nfw\PYZus{}profiles} \PY{o}{=} \PY{n}{np}\PY{o}{.}\PY{n}{load}\PY{p}{(}\PY{l+s+sa}{f}\PY{l+s+s1}{\PYZsq{}}\PY{l+s+si}{\PYZob{}}\PY{n}{obs\PYZus{}output\PYZus{}dir}\PY{l+s+si}{\PYZcb{}}\PY{l+s+s1}{/drawn\PYZus{}nfw\PYZus{}profiles.npy}\PY{l+s+s1}{\PYZsq{}}\PY{p}{)}


\PY{c+c1}{\PYZsh{} drawn\PYZus{}mc\PYZus{}pairs scatter plot}
\PY{n}{plt}\PY{o}{.}\PY{n}{scatter}\PY{p}{(}\PY{n}{drawn\PYZus{}mc\PYZus{}pairs}\PY{p}{[}\PY{p}{:}\PY{p}{,}\PY{l+m+mi}{0}\PY{p}{]}\PY{p}{,} \PY{n}{drawn\PYZus{}mc\PYZus{}pairs}\PY{p}{[}\PY{p}{:}\PY{p}{,}\PY{l+m+mi}{1}\PY{p}{]}\PY{p}{,}\PY{n}{s}\PY{o}{=}\PY{l+m+mi}{50}\PY{p}{)}
\PY{n}{plt}\PY{o}{.}\PY{n}{scatter}\PY{p}{(}
    \PY{n}{np}\PY{o}{.}\PY{n}{mean}\PY{p}{(}\PY{n}{drawn\PYZus{}mc\PYZus{}pairs}\PY{p}{[}\PY{p}{:}\PY{p}{,} \PY{l+m+mi}{0}\PY{p}{]}\PY{p}{)}\PY{p}{,}
    \PY{n}{np}\PY{o}{.}\PY{n}{mean}\PY{p}{(}\PY{n}{drawn\PYZus{}mc\PYZus{}pairs}\PY{p}{[}\PY{p}{:}\PY{p}{,} \PY{l+m+mi}{1}\PY{p}{]}\PY{p}{)}\PY{p}{,}
    \PY{n}{s}\PY{o}{=}\PY{l+m+mi}{100}\PY{p}{,}
    \PY{n}{marker}\PY{o}{=}\PY{l+s+s1}{\PYZsq{}}\PY{l+s+s1}{\PYZca{}}\PY{l+s+s1}{\PYZsq{}}\PY{p}{,}
    \PY{n}{label}\PY{o}{=}\PY{l+s+s1}{\PYZsq{}}\PY{l+s+s1}{mean m\PYZhy{}c pair}\PY{l+s+s1}{\PYZsq{}}\PY{p}{,}
\PY{p}{)}
\PY{n}{plt}\PY{o}{.}\PY{n}{legend}\PY{p}{(}\PY{p}{)}
\PY{n}{plt}\PY{o}{.}\PY{n}{xlabel}\PY{p}{(}\PY{l+s+s1}{\PYZsq{}}\PY{l+s+s1}{log\PYZdl{}\PYZus{}}\PY{l+s+si}{\PYZob{}10\PYZcb{}}\PY{l+s+s1}{\PYZdl{}M [M\PYZdl{}\PYZus{}}\PY{l+s+s1}{\PYZbs{}}\PY{l+s+s1}{odot\PYZdl{}]}\PY{l+s+s1}{\PYZsq{}}\PY{p}{,} \PY{n}{fontsize}\PY{o}{=}\PY{l+s+s1}{\PYZsq{}}\PY{l+s+s1}{xx\PYZhy{}large}\PY{l+s+s1}{\PYZsq{}}\PY{p}{)}
\PY{n}{plt}\PY{o}{.}\PY{n}{ylabel}\PY{p}{(}\PY{l+s+s1}{\PYZsq{}}\PY{l+s+s1}{Concentration}\PY{l+s+s1}{\PYZsq{}}\PY{p}{,} \PY{n}{fontsize}\PY{o}{=}\PY{l+s+s1}{\PYZsq{}}\PY{l+s+s1}{xx\PYZhy{}large}\PY{l+s+s1}{\PYZsq{}}\PY{p}{)}
\PY{n}{plt}\PY{o}{.}\PY{n}{title}\PY{p}{(}\PY{l+s+s1}{\PYZsq{}}\PY{l+s+s1}{Drawn mc\PYZus{}pairs}\PY{l+s+s1}{\PYZsq{}}\PY{p}{)}
\PY{n}{plt}\PY{o}{.}\PY{n}{show}\PY{p}{(}\PY{p}{)}
\end{Verbatim}
\end{tcolorbox}

    \begin{center}
    \adjustimage{max size={0.9\linewidth}{0.9\paperheight}}{tutorial_notebook_files/tutorial_notebook_16_0.png}
    \end{center}
    { \hspace*{\fill} \\}
    
    Next, we can look at the drawn\_nfw\_profiles which we generated from
the above drawn\_mc\_pairs.

Here, each individual NFW profile is noisy so we can perform stacking to
marginalize over this noise. One approach to stacking is that we can
first stack the observables and perform inference on that median NFW
profile. We call this approach ``join-then-fit'' because we first join
the observables and then fit on that ``joined'' observable.

    \begin{tcolorbox}[breakable, size=fbox, boxrule=1pt, pad at break*=1mm,colback=cellbackground, colframe=cellborder]
\prompt{In}{incolor}{23}{\boxspacing}
\begin{Verbatim}[commandchars=\\\{\}]
\PY{c+c1}{\PYZsh{} Expected shape: (NUM\PYZus{}OBS, num\PYZus{}radial\PYZus{}bins)}
\PY{n+nb}{print}\PY{p}{(}\PY{n}{np}\PY{o}{.}\PY{n}{shape}\PY{p}{(}\PY{n}{drawn\PYZus{}nfw\PYZus{}profiles}\PY{p}{)}\PY{p}{)}
\end{Verbatim}
\end{tcolorbox}

    \begin{Verbatim}[commandchars=\\\{\}]
(10, 30)
    \end{Verbatim}

    \begin{tcolorbox}[breakable, size=fbox, boxrule=1pt, pad at break*=1mm,colback=cellbackground, colframe=cellborder]
\prompt{In}{incolor}{123}{\boxspacing}
\begin{Verbatim}[commandchars=\\\{\}]
\PY{k+kn}{from} \PY{n+nn}{colossus}\PY{n+nn}{.}\PY{n+nn}{cosmology} \PY{k+kn}{import} \PY{n}{cosmology}
\PY{n}{cosmo} \PY{o}{=} \PY{n}{cosmology}\PY{o}{.}\PY{n}{setCosmology}\PY{p}{(}\PY{l+s+s1}{\PYZsq{}}\PY{l+s+s1}{planck18}\PY{l+s+s1}{\PYZsq{}}\PY{p}{)}

\PY{n}{rbins} \PY{o}{=} \PY{l+m+mi}{10}\PY{o}{*}\PY{o}{*}\PY{n}{np}\PY{o}{.}\PY{n}{arange}\PY{p}{(}\PY{l+m+mi}{0}\PY{p}{,} \PY{n}{sim\PYZus{}config}\PY{p}{[}\PY{l+s+s1}{\PYZsq{}}\PY{l+s+s1}{num\PYZus{}radial\PYZus{}bins}\PY{l+s+s1}{\PYZsq{}}\PY{p}{]} \PY{o}{/} \PY{l+m+mi}{10}\PY{p}{,} \PY{l+m+mf}{0.1}\PY{p}{)}
\PY{n}{plt}\PY{o}{.}\PY{n}{figure}\PY{p}{(}\PY{p}{)}
\PY{n}{plt}\PY{o}{.}\PY{n}{loglog}\PY{p}{(}\PY{p}{)}
\PY{n}{plt}\PY{o}{.}\PY{n}{xlabel}\PY{p}{(}\PY{l+s+s1}{\PYZsq{}}\PY{l+s+s1}{radius [kpc/h]}\PY{l+s+s1}{\PYZsq{}}\PY{p}{,} \PY{n}{fontsize}\PY{o}{=}\PY{l+s+s1}{\PYZsq{}}\PY{l+s+s1}{xx\PYZhy{}large}\PY{l+s+s1}{\PYZsq{}}\PY{p}{)}
\PY{n}{plt}\PY{o}{.}\PY{n}{ylabel}\PY{p}{(}\PY{l+s+s1}{\PYZsq{}}\PY{l+s+s1}{\PYZdl{} }\PY{l+s+s1}{\PYZbs{}}\PY{l+s+s1}{Delta }\PY{l+s+s1}{\PYZbs{}}\PY{l+s+s1}{Sigma\PYZdl{} [\PYZdl{}M\PYZus{}}\PY{l+s+s1}{\PYZbs{}}\PY{l+s+s1}{odot h / kpc\PYZca{}2\PYZdl{}]}\PY{l+s+s1}{\PYZsq{}}\PY{p}{,} \PY{n}{fontsize}\PY{o}{=}\PY{l+s+s1}{\PYZsq{}}\PY{l+s+s1}{xx\PYZhy{}large}\PY{l+s+s1}{\PYZsq{}}\PY{p}{)}

\PY{c+c1}{\PYZsh{} Plot each individual observation}
\PY{k}{for} \PY{n}{nfw\PYZus{}profile} \PY{o+ow}{in} \PY{n}{drawn\PYZus{}nfw\PYZus{}profiles}\PY{p}{:}
    \PY{n}{plt}\PY{o}{.}\PY{n}{plot}\PY{p}{(}\PY{n}{rbins}\PY{p}{,} \PY{n}{nfw\PYZus{}profile}\PY{p}{,} \PY{l+s+s1}{\PYZsq{}}\PY{l+s+s1}{\PYZhy{}}\PY{l+s+s1}{\PYZsq{}}\PY{p}{,} \PY{n}{alpha}\PY{o}{=}\PY{l+m+mf}{0.3}\PY{p}{,} \PY{n}{color}\PY{o}{=}\PY{l+s+s1}{\PYZsq{}}\PY{l+s+s1}{gray}\PY{l+s+s1}{\PYZsq{}}\PY{p}{)}

\PY{c+c1}{\PYZsh{} Calculate error bars}
\PY{n}{sigmas} \PY{o}{=} \PY{n}{np}\PY{o}{.}\PY{n}{std}\PY{p}{(}\PY{n}{drawn\PYZus{}nfw\PYZus{}profiles}\PY{p}{,} \PY{n}{axis}\PY{o}{=}\PY{l+m+mi}{0}\PY{p}{)}
\PY{n}{log\PYZus{}sigmas} \PY{o}{=} \PY{n}{np}\PY{o}{.}\PY{n}{std}\PY{p}{(}\PY{n}{np}\PY{o}{.}\PY{n}{log10}\PY{p}{(}\PY{n}{drawn\PYZus{}nfw\PYZus{}profiles}\PY{p}{)}\PY{p}{,} \PY{n}{axis}\PY{o}{=}\PY{l+m+mi}{0}\PY{p}{)}
\PY{n}{upper\PYZus{}error} \PY{o}{=} \PY{n}{np}\PY{o}{.}\PY{n}{exp}\PY{p}{(}
    \PY{n}{np}\PY{o}{.}\PY{n}{log}\PY{p}{(}\PY{n}{np}\PY{o}{.}\PY{n}{median}\PY{p}{(}\PY{n}{drawn\PYZus{}nfw\PYZus{}profiles}\PY{p}{,} \PY{n}{axis}\PY{o}{=}\PY{l+m+mi}{0}\PY{p}{)}\PY{p}{)} \PY{o}{+} \PY{n}{log\PYZus{}sigmas}\PY{p}{)} \PY{o}{\PYZhy{}} \PY{n}{np}\PY{o}{.}\PY{n}{median}\PY{p}{(}
        \PY{n}{drawn\PYZus{}nfw\PYZus{}profiles}\PY{p}{,} \PY{n}{axis}\PY{o}{=}\PY{l+m+mi}{0}\PY{p}{)}
\PY{n}{lower\PYZus{}error} \PY{o}{=} \PY{n}{np}\PY{o}{.}\PY{n}{median}\PY{p}{(}\PY{n}{drawn\PYZus{}nfw\PYZus{}profiles}\PY{p}{,} \PY{n}{axis}\PY{o}{=}\PY{l+m+mi}{0}\PY{p}{)} \PY{o}{\PYZhy{}} \PY{n}{np}\PY{o}{.}\PY{n}{exp}\PY{p}{(}
    \PY{n}{np}\PY{o}{.}\PY{n}{log}\PY{p}{(}\PY{n}{np}\PY{o}{.}\PY{n}{median}\PY{p}{(}\PY{n}{drawn\PYZus{}nfw\PYZus{}profiles}\PY{p}{,} \PY{n}{axis}\PY{o}{=}\PY{l+m+mi}{0}\PY{p}{)}\PY{p}{)} \PY{o}{\PYZhy{}} \PY{n}{log\PYZus{}sigmas}\PY{p}{)}
\PY{n}{yerr} \PY{o}{=} \PY{p}{[}\PY{n}{lower\PYZus{}error}\PY{p}{,} \PY{n}{upper\PYZus{}error}\PY{p}{]}

\PY{c+c1}{\PYZsh{} Plot median observation with error bars}
\PY{n}{plt}\PY{o}{.}\PY{n}{errorbar}\PY{p}{(}
    \PY{n}{rbins}\PY{p}{,}
    \PY{n}{np}\PY{o}{.}\PY{n}{median}\PY{p}{(}\PY{n}{drawn\PYZus{}nfw\PYZus{}profiles}\PY{p}{,} \PY{n}{axis}\PY{o}{=}\PY{l+m+mi}{0}\PY{p}{)}\PY{p}{,}
    \PY{n}{yerr}\PY{o}{=}\PY{n}{yerr}\PY{p}{,}
    \PY{n}{fmt}\PY{o}{=}\PY{l+s+s1}{\PYZsq{}}\PY{l+s+s1}{r}\PY{l+s+s1}{\PYZsq{}}\PY{p}{,}
    \PY{n}{capsize}\PY{o}{=}\PY{l+m+mf}{3.}\PY{p}{,}
    \PY{n}{label}\PY{o}{=}\PY{l+s+sa}{f}\PY{l+s+s1}{\PYZsq{}}\PY{l+s+s1}{Median Drawn NFW, }\PY{l+s+si}{\PYZob{}}\PY{n}{obs\PYZus{}config}\PY{p}{[}\PY{l+s+s2}{\PYZdq{}}\PY{l+s+s2}{min\PYZus{}richness}\PY{l+s+s2}{\PYZdq{}}\PY{p}{]}\PY{l+s+si}{\PYZcb{}}\PY{l+s+s1}{ \PYZlt{} \PYZdl{}}\PY{l+s+s1}{\PYZbs{}}\PY{l+s+s1}{lambda\PYZdl{} \PYZlt{} }\PY{l+s+si}{\PYZob{}}\PY{n}{obs\PYZus{}config}\PY{p}{[}\PY{l+s+s2}{\PYZdq{}}\PY{l+s+s2}{max\PYZus{}richness}\PY{l+s+s2}{\PYZdq{}}\PY{p}{]}\PY{l+s+si}{\PYZcb{}}\PY{l+s+s1}{\PYZsq{}}
\PY{p}{)}
\PY{n}{plt}\PY{o}{.}\PY{n}{legend}\PY{p}{(}\PY{p}{)}\PY{p}{;}
\end{Verbatim}
\end{tcolorbox}

    \begin{center}
    \adjustimage{max size={0.9\linewidth}{0.9\paperheight}}{tutorial_notebook_files/tutorial_notebook_19_0.png}
    \end{center}
    { \hspace*{\fill} \\}
    
    \hypertarget{inference}{%
\section{Inference}\label{inference}}

Now we'll need to infer the relationship between the input
log10mass-concentration pairs and the output NFW profiles. MCMC will
perform this inference using an analytic likelihood function. SBI will
perform inference this with a neural network inferrer that has been
trained on the simulated data points and their corresponding NFW
profiles. Both of these inference techniques will have the observations
applied to them, which can then show us the posterior distribution.

\hypertarget{configuration}{%
\subsection{Configuration}\label{configuration}}

Similarly to what we did before, we'll generate an infer\_config in our
directory of interest and look at an example of what these values could
be.

    \begin{tcolorbox}[breakable, size=fbox, boxrule=1pt, pad at break*=1mm,colback=cellbackground, colframe=cellborder]
\prompt{In}{incolor}{25}{\boxspacing}
\begin{Verbatim}[commandchars=\\\{\}]
\PY{c+c1}{\PYZsh{} Command line arguments for generating simulations}
\PY{n}{INFER\PYZus{}ID} \PY{o}{=} \PY{l+s+s1}{\PYZsq{}}\PY{l+s+s1}{infer\PYZus{}z1}\PY{l+s+s1}{\PYZsq{}}

\PY{k}{with} \PY{n+nb}{open}\PY{p}{(}\PY{l+s+sa}{f}\PY{l+s+s1}{\PYZsq{}}\PY{l+s+s1}{configs/inference/}\PY{l+s+si}{\PYZob{}}\PY{n}{INFER\PYZus{}ID}\PY{l+s+si}{\PYZcb{}}\PY{l+s+s1}{.json}\PY{l+s+s1}{\PYZsq{}}\PY{p}{,} \PY{l+s+s1}{\PYZsq{}}\PY{l+s+s1}{r}\PY{l+s+s1}{\PYZsq{}}\PY{p}{)} \PY{k}{as} \PY{n}{f}\PY{p}{:}
    \PY{n}{infer\PYZus{}config} \PY{o}{=} \PY{n}{json}\PY{o}{.}\PY{n}{load}\PY{p}{(}\PY{n}{f}\PY{p}{)}
    
\PY{n+nb}{print}\PY{p}{(}\PY{n}{json}\PY{o}{.}\PY{n}{dumps}\PY{p}{(}\PY{n}{infer\PYZus{}config}\PY{p}{,} \PY{n}{indent}\PY{o}{=}\PY{l+m+mi}{2}\PY{p}{)}\PY{p}{)}
\end{Verbatim}
\end{tcolorbox}

    \begin{Verbatim}[commandchars=\\\{\}]
\{
  "priors": \{
    "min\_log10mass": 10,
    "max\_log10mass": 20,
    "mc\_scatter": 0.1,
    "mc\_relation": "child18",
    "min\_concentration": 0,
    "max\_concentration": 10,
    "min\_z": 0.2,
    "max\_z": 0.35
  \}
\}
    \end{Verbatim}

    \hypertarget{run-inference}{%
\subsection{Run Inference}\label{run-inference}}

To actually run inference, we do the following: 1) train the SBI
inferrer, which (despite the misleading nomenclature of ``posterior'')
is actually just the trained inferrer (i.e.~it will become the posterior
later when we apply it to observations) 2) run inference 1) with MCMC
(using the likelihood function and observations) 2) with SBI (using
simulations, observations, and trained inferrer)

The reason we split these two steps is so we can train an SBI inferrer
and apply it to many different sets of observations without redundant
computing.

    \begin{tcolorbox}[breakable, size=fbox, boxrule=1pt, pad at break*=1mm,colback=cellbackground, colframe=cellborder]
\prompt{In}{incolor}{27}{\boxspacing}
\begin{Verbatim}[commandchars=\\\{\}]
\PY{o}{\PYZpc{}}\PY{k}{run} scripts/gen\PYZus{}posterior.py \PYZhy{}\PYZhy{}sim\PYZus{}id \PYZob{}SIM\PYZus{}ID\PYZcb{} \PYZhy{}\PYZhy{}infer\PYZus{}id \PYZob{}INFER\PYZus{}ID\PYZcb{} \PYZhy{}\PYZhy{}num\PYZus{}sims \PYZob{}NUM\PYZus{}SIMS\PYZcb{} \PYZhy{}\PYZhy{}regenerate
\end{Verbatim}
\end{tcolorbox}

    \begin{Verbatim}[commandchars=\\\{\}]
Overwriting existing posterior because of --regenerate flag
 Neural network successfully converged after 48 epochs.
    \end{Verbatim}

    This is a pretty time intensive step (\textasciitilde5 minutes). If you
want to stretch your legs and/or refill your cup of coffee, this is the
perfect time to do so. We are using SBI and MCMC to perform inference
here on NUM\_OBS + 1 observations: one for each of the observations and
one for the median of these observations.

While you're enjoying your coffee and running this step, we can discuss
the second method of stacking. We already mentioned ``join-then-fit''
where we stack observations first and perform inference on that joined
observation. The second stacking technique we use is called
``fit-then-join'' where we perform inference on each individual
observation and then stack the posteriors afterwards. We discuss the
pros and cons of each technique for each SBI and MCMC in the paper, but
at a high level, we find that join-then-fit results in tighter
constraints, but has the potential to include some bias. Fit-then-join,
by contrast, is often less constrained in parameter space, but is less
vulnerable to bias.

    \begin{tcolorbox}[breakable, size=fbox, boxrule=1pt, pad at break*=1mm,colback=cellbackground, colframe=cellborder]
\prompt{In}{incolor}{29}{\boxspacing}
\begin{Verbatim}[commandchars=\\\{\}]
\PY{o}{\PYZpc{}}\PY{k}{run} scripts/run\PYZus{}inference.py \PYZhy{}\PYZhy{}sim\PYZus{}id \PYZob{}SIM\PYZus{}ID\PYZcb{} \PYZhy{}\PYZhy{}infer\PYZus{}id \PYZob{}INFER\PYZus{}ID\PYZcb{} \PYZhy{}\PYZhy{}obs\PYZus{}id \PYZob{}OBS\PYZus{}ID\PYZcb{} \PYZhy{}\PYZhy{}num\PYZus{}sims \PYZob{}NUM\PYZus{}SIMS\PYZcb{} \PYZhy{}\PYZhy{}num\PYZus{}obs \PYZob{}NUM\PYZus{}OBS\PYZcb{} \PYZhy{}\PYZhy{}regenerate
\PY{n+nb}{print}\PY{p}{(}\PY{l+s+s2}{\PYZdq{}}\PY{l+s+s2}{Completed running inference}\PY{l+s+s2}{\PYZdq{}}\PY{p}{)}
\end{Verbatim}
\end{tcolorbox}

    \begin{Verbatim}[commandchars=\\\{\}]
Overwriting existing inference outputs because of --regenerate flag
    \end{Verbatim}

    
    \begin{Verbatim}[commandchars=\\\{\}]
Drawing 10000 posterior samples:   0\%|          | 0/10000 [00:00<?, ?it/s]
    \end{Verbatim}

    
    \begin{Verbatim}[commandchars=\\\{\}]
/Users/akumgill/Library/Python/3.9/lib/python/site-
packages/pyknos/mdn/mdn.py:283: UserWarning: torch.triangular\_solve is
deprecated in favor of torch.linalg.solve\_triangularand will be removed in a
future PyTorch release.
torch.linalg.solve\_triangular has its arguments reversed and does not return a
copy of one of the inputs.
X = torch.triangular\_solve(B, A).solution
should be replaced with
X = torch.linalg.solve\_triangular(A, B). (Triggered internally at /Users/runner/
work/pytorch/pytorch/pytorch/aten/src/ATen/native/BatchLinearAlgebra.cpp:2197.)
  zero\_mean\_samples, \_ = torch.triangular\_solve(
    \end{Verbatim}

    
    \begin{Verbatim}[commandchars=\\\{\}]
Drawing 10000 posterior samples:   0\%|          | 0/10000 [00:00<?, ?it/s]
    \end{Verbatim}

    
    
    \begin{Verbatim}[commandchars=\\\{\}]
Drawing 10000 posterior samples:   0\%|          | 0/10000 [00:00<?, ?it/s]
    \end{Verbatim}

    
    
    \begin{Verbatim}[commandchars=\\\{\}]
Drawing 10000 posterior samples:   0\%|          | 0/10000 [00:00<?, ?it/s]
    \end{Verbatim}

    
    
    \begin{Verbatim}[commandchars=\\\{\}]
Drawing 10000 posterior samples:   0\%|          | 0/10000 [00:00<?, ?it/s]
    \end{Verbatim}

    
    
    \begin{Verbatim}[commandchars=\\\{\}]
Drawing 10000 posterior samples:   0\%|          | 0/10000 [00:00<?, ?it/s]
    \end{Verbatim}

    
    
    \begin{Verbatim}[commandchars=\\\{\}]
Drawing 10000 posterior samples:   0\%|          | 0/10000 [00:00<?, ?it/s]
    \end{Verbatim}

    
    
    \begin{Verbatim}[commandchars=\\\{\}]
Drawing 10000 posterior samples:   0\%|          | 0/10000 [00:00<?, ?it/s]
    \end{Verbatim}

    
    
    \begin{Verbatim}[commandchars=\\\{\}]
Drawing 10000 posterior samples:   0\%|          | 0/10000 [00:00<?, ?it/s]
    \end{Verbatim}

    
    
    \begin{Verbatim}[commandchars=\\\{\}]
Drawing 10000 posterior samples:   0\%|          | 0/10000 [00:00<?, ?it/s]
    \end{Verbatim}

    
    
    \begin{Verbatim}[commandchars=\\\{\}]
Drawing 10000 posterior samples:   0\%|          | 0/10000 [00:00<?, ?it/s]
    \end{Verbatim}

    
    \begin{Verbatim}[commandchars=\\\{\}]
\#\# burning in {\ldots}
    \end{Verbatim}

    \begin{Verbatim}[commandchars=\\\{\}]
/Users/akumgill/Documents/GitHub/CL-
SBI/weaklensclustersbi/inference/mcmcutils.py:24: RuntimeWarning: divide by zero
encountered in log
  return np.log(c\_norm / 20)
  0\%|                                           | 0/100 [00:00<?,
?it/s]/Users/akumgill/Library/Python/3.9/lib/python/site-
packages/colossus/halo/concentration.py:442: UserWarning: Some masses or
redshifts are outside the validity of the concentration model.
  warnings.warn('Some masses or redshifts are outside the validity of the
concentration model.')
/Users/akumgill/Library/Python/3.9/lib/python/site-
packages/emcee/moves/red\_blue.py:99: RuntimeWarning: invalid value encountered
in scalar subtract
  lnpdiff = f + nlp - state.log\_prob[j]
100\%|█████████████████████████████████| 100/100 [00:14<00:00,  6.85it/s]
    \end{Verbatim}

    \begin{Verbatim}[commandchars=\\\{\}]
\#\# running the full chain {\ldots}
    \end{Verbatim}

    \begin{Verbatim}[commandchars=\\\{\}]
100\%|█████████████████████████████████| 200/200 [00:30<00:00,  6.58it/s]
    \end{Verbatim}

    \begin{Verbatim}[commandchars=\\\{\}]
\#\# burning in {\ldots}
    \end{Verbatim}

    \begin{Verbatim}[commandchars=\\\{\}]
100\%|█████████████████████████████████| 100/100 [00:14<00:00,  6.73it/s]
    \end{Verbatim}

    \begin{Verbatim}[commandchars=\\\{\}]
\#\# running the full chain {\ldots}
    \end{Verbatim}

    \begin{Verbatim}[commandchars=\\\{\}]
100\%|█████████████████████████████████| 200/200 [00:31<00:00,  6.31it/s]
    \end{Verbatim}

    \begin{Verbatim}[commandchars=\\\{\}]
\#\# burning in {\ldots}
    \end{Verbatim}

    \begin{Verbatim}[commandchars=\\\{\}]
100\%|█████████████████████████████████| 100/100 [00:14<00:00,  6.84it/s]
    \end{Verbatim}

    \begin{Verbatim}[commandchars=\\\{\}]
\#\# running the full chain {\ldots}
    \end{Verbatim}

    \begin{Verbatim}[commandchars=\\\{\}]
100\%|█████████████████████████████████| 200/200 [00:30<00:00,  6.50it/s]
    \end{Verbatim}

    \begin{Verbatim}[commandchars=\\\{\}]
\#\# burning in {\ldots}
    \end{Verbatim}

    \begin{Verbatim}[commandchars=\\\{\}]
100\%|█████████████████████████████████| 100/100 [00:14<00:00,  6.85it/s]
    \end{Verbatim}

    \begin{Verbatim}[commandchars=\\\{\}]
\#\# running the full chain {\ldots}
    \end{Verbatim}

    \begin{Verbatim}[commandchars=\\\{\}]
100\%|█████████████████████████████████| 200/200 [00:31<00:00,  6.31it/s]
    \end{Verbatim}

    \begin{Verbatim}[commandchars=\\\{\}]
\#\# burning in {\ldots}
    \end{Verbatim}

    \begin{Verbatim}[commandchars=\\\{\}]
100\%|█████████████████████████████████| 100/100 [00:14<00:00,  6.80it/s]
    \end{Verbatim}

    \begin{Verbatim}[commandchars=\\\{\}]
\#\# running the full chain {\ldots}
    \end{Verbatim}

    \begin{Verbatim}[commandchars=\\\{\}]
100\%|█████████████████████████████████| 200/200 [00:31<00:00,  6.37it/s]
    \end{Verbatim}

    \begin{Verbatim}[commandchars=\\\{\}]
\#\# burning in {\ldots}
    \end{Verbatim}

    \begin{Verbatim}[commandchars=\\\{\}]
100\%|█████████████████████████████████| 100/100 [00:15<00:00,  6.56it/s]
    \end{Verbatim}

    \begin{Verbatim}[commandchars=\\\{\}]
\#\# running the full chain {\ldots}
    \end{Verbatim}

    \begin{Verbatim}[commandchars=\\\{\}]
100\%|█████████████████████████████████| 200/200 [00:31<00:00,  6.33it/s]
    \end{Verbatim}

    \begin{Verbatim}[commandchars=\\\{\}]
\#\# burning in {\ldots}
    \end{Verbatim}

    \begin{Verbatim}[commandchars=\\\{\}]
100\%|█████████████████████████████████| 100/100 [00:15<00:00,  6.29it/s]
    \end{Verbatim}

    \begin{Verbatim}[commandchars=\\\{\}]
\#\# running the full chain {\ldots}
    \end{Verbatim}

    \begin{Verbatim}[commandchars=\\\{\}]
100\%|█████████████████████████████████| 200/200 [00:32<00:00,  6.12it/s]
    \end{Verbatim}

    \begin{Verbatim}[commandchars=\\\{\}]
\#\# burning in {\ldots}
    \end{Verbatim}

    \begin{Verbatim}[commandchars=\\\{\}]
100\%|█████████████████████████████████| 100/100 [00:15<00:00,  6.62it/s]
    \end{Verbatim}

    \begin{Verbatim}[commandchars=\\\{\}]
\#\# running the full chain {\ldots}
    \end{Verbatim}

    \begin{Verbatim}[commandchars=\\\{\}]
100\%|█████████████████████████████████| 200/200 [00:31<00:00,  6.39it/s]
    \end{Verbatim}

    \begin{Verbatim}[commandchars=\\\{\}]
\#\# burning in {\ldots}
    \end{Verbatim}

    \begin{Verbatim}[commandchars=\\\{\}]
100\%|█████████████████████████████████| 100/100 [00:14<00:00,  6.75it/s]
    \end{Verbatim}

    \begin{Verbatim}[commandchars=\\\{\}]
\#\# running the full chain {\ldots}
    \end{Verbatim}

    \begin{Verbatim}[commandchars=\\\{\}]
100\%|█████████████████████████████████| 200/200 [00:30<00:00,  6.47it/s]
    \end{Verbatim}

    \begin{Verbatim}[commandchars=\\\{\}]
\#\# burning in {\ldots}
    \end{Verbatim}

    \begin{Verbatim}[commandchars=\\\{\}]
100\%|█████████████████████████████████| 100/100 [00:14<00:00,  6.75it/s]
    \end{Verbatim}

    \begin{Verbatim}[commandchars=\\\{\}]
\#\# running the full chain {\ldots}
    \end{Verbatim}

    \begin{Verbatim}[commandchars=\\\{\}]
100\%|█████████████████████████████████| 200/200 [00:30<00:00,  6.56it/s]
    \end{Verbatim}

    \begin{Verbatim}[commandchars=\\\{\}]
\#\# burning in {\ldots}
    \end{Verbatim}

    \begin{Verbatim}[commandchars=\\\{\}]
100\%|█████████████████████████████████| 100/100 [00:15<00:00,  6.66it/s]
    \end{Verbatim}

    \begin{Verbatim}[commandchars=\\\{\}]
\#\# running the full chain {\ldots}
    \end{Verbatim}

    \begin{Verbatim}[commandchars=\\\{\}]
100\%|█████████████████████████████████| 200/200 [00:30<00:00,  6.49it/s]
    \end{Verbatim}

    \begin{Verbatim}[commandchars=\\\{\}]
Completed running inference
    \end{Verbatim}

    \begin{Verbatim}[commandchars=\\\{\}]

    \end{Verbatim}

    \hypertarget{plotting}{%
\subsection{Plotting}\label{plotting}}

Now we can plot the confusogram/corner plots/banana plots.

    \begin{tcolorbox}[breakable, size=fbox, boxrule=1pt, pad at break*=1mm,colback=cellbackground, colframe=cellborder]
\prompt{In}{incolor}{56}{\boxspacing}
\begin{Verbatim}[commandchars=\\\{\}]
\PY{c+c1}{\PYZsh{} Turning off interactive plotting so we can be more intentional}
\PY{c+c1}{\PYZsh{} plt.ioff();}
\PY{o}{\PYZpc{}}\PY{k}{run} scripts/plot\PYZus{}chains.py \PYZhy{}\PYZhy{}sim\PYZus{}id \PYZob{}SIM\PYZus{}ID\PYZcb{} \PYZhy{}\PYZhy{}infer\PYZus{}id \PYZob{}INFER\PYZus{}ID\PYZcb{} \PYZhy{}\PYZhy{}obs\PYZus{}id \PYZob{}OBS\PYZus{}ID\PYZcb{} \PYZhy{}\PYZhy{}num\PYZus{}sims \PYZob{}NUM\PYZus{}SIMS\PYZcb{} \PYZhy{}\PYZhy{}num\PYZus{}obs \PYZob{}NUM\PYZus{}OBS\PYZcb{} \PYZhy{}\PYZhy{}regenerate
\PY{n}{plt}\PY{o}{.}\PY{n}{close}\PY{p}{(}\PY{p}{)}
\PY{c+c1}{\PYZsh{} plt.ion();}
\end{Verbatim}
\end{tcolorbox}

    \begin{Verbatim}[commandchars=\\\{\}]
Overwriting existing plots because of --regenerate flag
    \end{Verbatim}

    As before, let's first see what plots were generated. We see mostly
files like \texttt{\{inference\}\_\{plotting\}.\{filetype\}}
(e.g.~\texttt{mcmc\_cc.pdf}). There is also a file called
\texttt{mcmc\_sbi\_cc.\{filetype\}} which shows both MCMC and SBI
inference together, but this can be tricky to parse. We also see a
\texttt{diagnostics/} folder which we will look into shortly.

    \begin{tcolorbox}[breakable, size=fbox, boxrule=1pt, pad at break*=1mm,colback=cellbackground, colframe=cellborder]
\prompt{In}{incolor}{57}{\boxspacing}
\begin{Verbatim}[commandchars=\\\{\}]
\PY{c+c1}{\PYZsh{} Confirming the output files were generated correctly}
\PY{n}{plots\PYZus{}output\PYZus{}dir} \PY{o}{=} \PY{l+s+sa}{f}\PY{l+s+s1}{\PYZsq{}}\PY{l+s+s1}{outputs/plots/}\PY{l+s+si}{\PYZob{}}\PY{n}{SIM\PYZus{}ID}\PY{l+s+si}{\PYZcb{}}\PY{l+s+s1}{.}\PY{l+s+si}{\PYZob{}}\PY{n}{INFER\PYZus{}ID}\PY{l+s+si}{\PYZcb{}}\PY{l+s+s1}{.}\PY{l+s+si}{\PYZob{}}\PY{n}{OBS\PYZus{}ID}\PY{l+s+si}{\PYZcb{}}\PY{l+s+s1}{.}\PY{l+s+si}{\PYZob{}}\PY{n}{NUM\PYZus{}SIMS}\PY{l+s+si}{\PYZcb{}}\PY{l+s+s1}{.}\PY{l+s+si}{\PYZob{}}\PY{n}{NUM\PYZus{}OBS}\PY{l+s+si}{\PYZcb{}}\PY{l+s+s1}{\PYZsq{}}
\PY{o}{\PYZpc{}}\PY{k}{ls} \PYZob{}plots\PYZus{}output\PYZus{}dir\PYZcb{}
\end{Verbatim}
\end{tcolorbox}

    \begin{Verbatim}[commandchars=\\\{\}]
\textcolor{ansi-cyan-intense}{\textbf{diagnostics}}/     mcmc\_gtc.pdf     mcmc\_sbi\_cc.png  sbi\_gtc.pdf
mcmc\_cc.pdf      mcmc\_gtc.png     sbi\_cc.pdf       sbi\_gtc.png
mcmc\_cc.png      mcmc\_sbi\_cc.pdf  sbi\_cc.png
    \end{Verbatim}

    \begin{tcolorbox}[breakable, size=fbox, boxrule=1pt, pad at break*=1mm,colback=cellbackground, colframe=cellborder]
\prompt{In}{incolor}{90}{\boxspacing}
\begin{Verbatim}[commandchars=\\\{\}]
\PY{n+nb}{print}\PY{p}{(}\PY{l+s+s1}{\PYZsq{}}\PY{l+s+s1}{MCMC corner plot using ChainConsumer}\PY{l+s+s1}{\PYZsq{}}\PY{p}{)}
\PY{n}{plt}\PY{o}{.}\PY{n}{figure}\PY{p}{(}\PY{n}{figsize}\PY{o}{=}\PY{p}{(}\PY{l+m+mi}{10}\PY{p}{,} \PY{l+m+mi}{8}\PY{p}{)}\PY{p}{)}
\PY{n}{plt}\PY{o}{.}\PY{n}{imshow}\PY{p}{(}\PY{n}{mpimg}\PY{o}{.}\PY{n}{imread}\PY{p}{(}\PY{l+s+sa}{f}\PY{l+s+s1}{\PYZsq{}}\PY{l+s+si}{\PYZob{}}\PY{n}{plots\PYZus{}output\PYZus{}dir}\PY{l+s+si}{\PYZcb{}}\PY{l+s+s1}{/mcmc\PYZus{}cc.png}\PY{l+s+s1}{\PYZsq{}}\PY{p}{)}\PY{p}{)}
\PY{n}{plt}\PY{o}{.}\PY{n}{axis}\PY{p}{(}\PY{l+s+s1}{\PYZsq{}}\PY{l+s+s1}{off}\PY{l+s+s1}{\PYZsq{}}\PY{p}{)}
\PY{n}{plt}\PY{o}{.}\PY{n}{show}\PY{p}{(}\PY{p}{)}

\PY{n+nb}{print}\PY{p}{(}\PY{l+s+s1}{\PYZsq{}}\PY{l+s+s1}{SBI corner plot using ChainConsumer}\PY{l+s+s1}{\PYZsq{}}\PY{p}{)}
\PY{n}{plt}\PY{o}{.}\PY{n}{figure}\PY{p}{(}\PY{n}{figsize}\PY{o}{=}\PY{p}{(}\PY{l+m+mi}{10}\PY{p}{,} \PY{l+m+mi}{8}\PY{p}{)}\PY{p}{)}
\PY{n}{plt}\PY{o}{.}\PY{n}{imshow}\PY{p}{(}\PY{n}{mpimg}\PY{o}{.}\PY{n}{imread}\PY{p}{(}\PY{l+s+sa}{f}\PY{l+s+s1}{\PYZsq{}}\PY{l+s+si}{\PYZob{}}\PY{n}{plots\PYZus{}output\PYZus{}dir}\PY{l+s+si}{\PYZcb{}}\PY{l+s+s1}{/sbi\PYZus{}cc.png}\PY{l+s+s1}{\PYZsq{}}\PY{p}{)}\PY{p}{)}
\PY{n}{plt}\PY{o}{.}\PY{n}{axis}\PY{p}{(}\PY{l+s+s1}{\PYZsq{}}\PY{l+s+s1}{off}\PY{l+s+s1}{\PYZsq{}}\PY{p}{)}
\PY{n}{plt}\PY{o}{.}\PY{n}{show}\PY{p}{(}\PY{p}{)}

\PY{n+nb}{print}\PY{p}{(}\PY{l+s+s1}{\PYZsq{}}\PY{l+s+s1}{MCMC and SBI corner plot using ChainConsumer. As you can see, this gets hard to parse, but is useful to show that the inference is consistent for this configuration across these two inference techniques.}\PY{l+s+s1}{\PYZsq{}}\PY{p}{)}
\PY{n}{plt}\PY{o}{.}\PY{n}{figure}\PY{p}{(}\PY{n}{figsize}\PY{o}{=}\PY{p}{(}\PY{l+m+mi}{10}\PY{p}{,} \PY{l+m+mi}{8}\PY{p}{)}\PY{p}{)}
\PY{n}{plt}\PY{o}{.}\PY{n}{imshow}\PY{p}{(}\PY{n}{mpimg}\PY{o}{.}\PY{n}{imread}\PY{p}{(}\PY{l+s+sa}{f}\PY{l+s+s1}{\PYZsq{}}\PY{l+s+si}{\PYZob{}}\PY{n}{plots\PYZus{}output\PYZus{}dir}\PY{l+s+si}{\PYZcb{}}\PY{l+s+s1}{/mcmc\PYZus{}sbi\PYZus{}cc.png}\PY{l+s+s1}{\PYZsq{}}\PY{p}{)}\PY{p}{)}
\PY{n}{plt}\PY{o}{.}\PY{n}{axis}\PY{p}{(}\PY{l+s+s1}{\PYZsq{}}\PY{l+s+s1}{off}\PY{l+s+s1}{\PYZsq{}}\PY{p}{)}
\PY{n}{plt}\PY{o}{.}\PY{n}{show}\PY{p}{(}\PY{p}{)}
\end{Verbatim}
\end{tcolorbox}

    \begin{Verbatim}[commandchars=\\\{\}]
MCMC corner plot using ChainConsumer
    \end{Verbatim}

    \begin{center}
    \adjustimage{max size={0.9\linewidth}{0.9\paperheight}}{tutorial_notebook_files/tutorial_notebook_30_1.png}
    \end{center}
    { \hspace*{\fill} \\}
    
    \begin{Verbatim}[commandchars=\\\{\}]
SBI corner plot using ChainConsumer
    \end{Verbatim}

    \begin{center}
    \adjustimage{max size={0.9\linewidth}{0.9\paperheight}}{tutorial_notebook_files/tutorial_notebook_30_3.png}
    \end{center}
    { \hspace*{\fill} \\}
    
    \begin{Verbatim}[commandchars=\\\{\}]
MCMC and SBI corner plot using ChainConsumer. As you can see, this gets hard to
parse, but is useful to show that the inference is consistent for this
configuration across these two inference techniques.
    \end{Verbatim}

    \begin{center}
    \adjustimage{max size={0.9\linewidth}{0.9\paperheight}}{tutorial_notebook_files/tutorial_notebook_30_5.png}
    \end{center}
    { \hspace*{\fill} \\}
    
    Finally, we can also look at the diagnostic plots. Here, we have some
useful files: * Drawn NFW profiles shows us whether the NFWs from our
inferred log10mass and concentration are comparable with the observed
ones * MCMC walkers will show us whether we have a sufficient number of
steps for MCMC to actually converge * \{inference\}\_ftj\_cc shows us
the inferred log10mass and concentration for each individual
observation. These posteriors are subsequently stacked in the
``fit-then-join'' stacking technique.

    \begin{tcolorbox}[breakable, size=fbox, boxrule=1pt, pad at break*=1mm,colback=cellbackground, colframe=cellborder]
\prompt{In}{incolor}{88}{\boxspacing}
\begin{Verbatim}[commandchars=\\\{\}]
\PY{o}{\PYZpc{}}\PY{k}{ls} \PYZob{}plots\PYZus{}output\PYZus{}dir\PYZcb{}/diagnostics
\end{Verbatim}
\end{tcolorbox}

    \begin{Verbatim}[commandchars=\\\{\}]
drawn\_nfw\_profiles.pdf     mcmc\_jtf\_mcmc\_walkers.pdf
drawn\_nfw\_profiles.png     mcmc\_jtf\_mcmc\_walkers.png
mcmc\_ftj\_cc.pdf            sbi\_ftj\_cc.pdf
mcmc\_ftj\_cc.png            sbi\_ftj\_cc.png
    \end{Verbatim}

    \begin{tcolorbox}[breakable, size=fbox, boxrule=1pt, pad at break*=1mm,colback=cellbackground, colframe=cellborder]
\prompt{In}{incolor}{98}{\boxspacing}
\begin{Verbatim}[commandchars=\\\{\}]
\PY{n+nb}{print}\PY{p}{(}\PY{l+s+s1}{\PYZsq{}}\PY{l+s+s1}{Inferred NFW profiles: We can see that the inferred parameters are consistent with the observations, which is exactly what we would expect.}\PY{l+s+s1}{\PYZsq{}}\PY{p}{)}
\PY{n}{plt}\PY{o}{.}\PY{n}{figure}\PY{p}{(}\PY{n}{figsize}\PY{o}{=}\PY{p}{(}\PY{l+m+mi}{10}\PY{p}{,} \PY{l+m+mi}{8}\PY{p}{)}\PY{p}{)}
\PY{n}{plt}\PY{o}{.}\PY{n}{imshow}\PY{p}{(}\PY{n}{mpimg}\PY{o}{.}\PY{n}{imread}\PY{p}{(}\PY{l+s+sa}{f}\PY{l+s+s1}{\PYZsq{}}\PY{l+s+si}{\PYZob{}}\PY{n}{plots\PYZus{}output\PYZus{}dir}\PY{l+s+si}{\PYZcb{}}\PY{l+s+s1}{/diagnostics/drawn\PYZus{}nfw\PYZus{}profiles.png}\PY{l+s+s1}{\PYZsq{}}\PY{p}{)}\PY{p}{)}
\PY{n}{plt}\PY{o}{.}\PY{n}{axis}\PY{p}{(}\PY{l+s+s1}{\PYZsq{}}\PY{l+s+s1}{off}\PY{l+s+s1}{\PYZsq{}}\PY{p}{)}
\PY{n}{plt}\PY{o}{.}\PY{n}{show}\PY{p}{(}\PY{p}{)}

\PY{n+nb}{print}\PY{p}{(}\PY{l+s+s1}{\PYZsq{}}\PY{l+s+s1}{MCMC Walkers: We see that after the burnout phase, our MCMC walkers have all converged on the relevant part of posterior space and are successfully exploring that region.}\PY{l+s+s1}{\PYZsq{}}\PY{p}{)}
\PY{n}{plt}\PY{o}{.}\PY{n}{figure}\PY{p}{(}\PY{n}{figsize}\PY{o}{=}\PY{p}{(}\PY{l+m+mi}{10}\PY{p}{,} \PY{l+m+mi}{8}\PY{p}{)}\PY{p}{)}
\PY{n}{plt}\PY{o}{.}\PY{n}{imshow}\PY{p}{(}\PY{n}{mpimg}\PY{o}{.}\PY{n}{imread}\PY{p}{(}\PY{l+s+sa}{f}\PY{l+s+s1}{\PYZsq{}}\PY{l+s+si}{\PYZob{}}\PY{n}{plots\PYZus{}output\PYZus{}dir}\PY{l+s+si}{\PYZcb{}}\PY{l+s+s1}{/diagnostics/mcmc\PYZus{}jtf\PYZus{}mcmc\PYZus{}walkers.png}\PY{l+s+s1}{\PYZsq{}}\PY{p}{)}\PY{p}{)}
\PY{n}{plt}\PY{o}{.}\PY{n}{axis}\PY{p}{(}\PY{l+s+s1}{\PYZsq{}}\PY{l+s+s1}{off}\PY{l+s+s1}{\PYZsq{}}\PY{p}{)}
\PY{n}{plt}\PY{o}{.}\PY{n}{show}\PY{p}{(}\PY{p}{)}

\PY{n+nb}{print}\PY{p}{(}\PY{l+s+s1}{\PYZsq{}}\PY{l+s+s1}{SBI contours for individual observations: These posteriors are subsequently stacked for fit\PYZhy{}then\PYZhy{}join (shown above).}\PY{l+s+s1}{\PYZsq{}}\PY{p}{)}
\PY{n}{plt}\PY{o}{.}\PY{n}{figure}\PY{p}{(}\PY{n}{figsize}\PY{o}{=}\PY{p}{(}\PY{l+m+mi}{10}\PY{p}{,} \PY{l+m+mi}{8}\PY{p}{)}\PY{p}{)}
\PY{n}{plt}\PY{o}{.}\PY{n}{imshow}\PY{p}{(}\PY{n}{mpimg}\PY{o}{.}\PY{n}{imread}\PY{p}{(}\PY{l+s+sa}{f}\PY{l+s+s1}{\PYZsq{}}\PY{l+s+si}{\PYZob{}}\PY{n}{plots\PYZus{}output\PYZus{}dir}\PY{l+s+si}{\PYZcb{}}\PY{l+s+s1}{/diagnostics/sbi\PYZus{}ftj\PYZus{}cc.png}\PY{l+s+s1}{\PYZsq{}}\PY{p}{)}\PY{p}{)}
\PY{n}{plt}\PY{o}{.}\PY{n}{axis}\PY{p}{(}\PY{l+s+s1}{\PYZsq{}}\PY{l+s+s1}{off}\PY{l+s+s1}{\PYZsq{}}\PY{p}{)}
\PY{n}{plt}\PY{o}{.}\PY{n}{show}\PY{p}{(}\PY{p}{)}
\end{Verbatim}
\end{tcolorbox}

    \begin{Verbatim}[commandchars=\\\{\}]
Inferred NFW profiles: We can see that the inferred parameters are consistent
with the observations, which is exactly what we would expect.
    \end{Verbatim}

    \begin{center}
    \adjustimage{max size={0.9\linewidth}{0.9\paperheight}}{tutorial_notebook_files/tutorial_notebook_33_1.png}
    \end{center}
    { \hspace*{\fill} \\}
    
    \begin{Verbatim}[commandchars=\\\{\}]
MCMC Walkers: We see that after the burnout phase, our MCMC walkers have all
converged on the relevant part of posterior space and are successfully exploring
that region.
    \end{Verbatim}

    \begin{center}
    \adjustimage{max size={0.9\linewidth}{0.9\paperheight}}{tutorial_notebook_files/tutorial_notebook_33_3.png}
    \end{center}
    { \hspace*{\fill} \\}
    
    \begin{Verbatim}[commandchars=\\\{\}]
SBI contours for individual observations: These posteriors are subsequently
stacked for fit-then-join (shown above).
    \end{Verbatim}

    \begin{center}
    \adjustimage{max size={0.9\linewidth}{0.9\paperheight}}{tutorial_notebook_files/tutorial_notebook_33_5.png}
    \end{center}
    { \hspace*{\fill} \\}
    

    % Add a bibliography block to the postdoc
    
    
    
\end{document}
